\documentclass[11pt,a4paper]{article}
\usepackage[T1]{fontenc}
\usepackage{isabelle,isabellesym}

\usepackage{amsmath}
\usepackage{amssymb}
\usepackage{stmaryrd}

% this should be the last package used
\usepackage{pdfsetup}

\urlstyle{rm}
\isabellestyle{it}

\begin{document}

\title{The Spectral Theory of Real Symmetric Matrices}
\author{Dmitriy Traytel}
\maketitle

\begin{abstract}
  The spectral theory of real symmetric matrices that HOL-Analysis stops
  short of, for anyone who constrains a matrix through its spectrum rather
  than through its entries.

  The spectral theorem itself is proved by the variational argument:
  a maximizer of the Rayleigh quotient on an invariant subspace is an
  eigenvector, and induction on dimension peels off an orthonormal eigenbasis
  of the whole space. On top of it, Ky Fan's theorem is developed in the form
  needed to constrain a matrix by its ordered eigenvalues without ever naming
  them individually: the sum of the $m$ largest eigenvalues is characterised
  as a supremum of traces against orthogonal projections, the individual
  ordered eigenvalues are recovered as differences of consecutive such sums,
  and both are shown Lipschitz in the matrix. The Poincar\'e separation
  inequality compares the spectrum of a matrix to that of its compression to
  a subspace.

  Nothing in the entry refers to any particular equation; it was extracted
  from a formalisation of the uniqueness half of a relative-arbitrage
  problem~\cite{LaiShkolnikovSoner}, where a lower bound on a sum of ordered
  eigenvalues is one of the two spectral constraints defining the problem's
  elliptic operator.
\end{abstract}

\tableofcontents

\parindent 0pt\parskip 0.5ex

\input{session}

\bibliographystyle{abbrv}
\bibliography{root}

\end{document}
